% =========================================================
% Completeness of Metric Spaces
% Source: volume-iii/analysis/metric-spaces/notes/notes-completeness.tex
%
% Dependencies:
%   notes-cauchy-sequences.tex — def:cauchy-sequence,
%                                prop:convergent-implies-cauchy,
%                                prop:cauchy-bounded,
%                                prop:cauchy-subsequence
%   sequential-characterizations.tex — sec:sequential-characterizations
%   open-closed-sets.tex       — def:open-closed-set,
%                                prop:closed-owns-boundary
% =========================================================

% ---------------------------------------------------------
% Toolkit
% ---------------------------------------------------------

\begin{tcolorbox}[colback=gray!6, colframe=gray!40, arc=2pt,
  left=6pt, right=6pt, top=4pt, bottom=4pt,
  title={\small\textbf{Completeness --- Quick Reference}},
  fonttitle=\small\bfseries]
\small
\begin{tabular}{l l l l}
\toprule
\textbf{Label} & \textbf{Statement} & \textbf{Proof method} & \textbf{Detail} \\
\midrule
D\ref{def:complete-metric-space}
  & $(X,d)$ complete $\iff$ every Cauchy sequence converges in $X$
  & Definition
  & \hyperref[def:complete-metric-space]{$\downarrow$} \\[4pt]
P\ref{prop:closed-subspace-complete}
  & Closed subspace of a complete space is complete
  & Sequential closedness
  & \hyperref[prop:closed-subspace-complete]{$\downarrow$} \\[4pt]
P\ref{prop:complete-subspace-closed}
  & Complete subspace of any metric space is closed
  & Uniqueness of limits
  & \hyperref[prop:complete-subspace-closed]{$\downarrow$} \\[4pt]
T\ref{thm:completion}
  & Every metric space has a completion, unique up to isometry
  & Co-Cauchy equivalence classes
  & \hyperref[thm:completion]{$\downarrow$} \\[4pt]
\bottomrule
\end{tabular}
\end{tcolorbox}

% ---------------------------------------------------------
\subsection{Complete Metric Spaces}
\label{sec:complete-metric-spaces}
% ---------------------------------------------------------

In $\mathbb{R}$, the Cauchy criterion is an \emph{if and only if}: a
sequence converges precisely when it is Cauchy.
This equivalence is a consequence of the completeness of $\mathbb{R}$
(the least upper bound property), and it fails in general metric spaces.
The notion of completeness captures exactly this property --- that no
Cauchy sequence is ``missing'' its limit.

\begin{tcolorbox}[colback=propbox, colframe=propborder, arc=2pt,
  left=6pt, right=6pt, top=4pt, bottom=4pt,
  title={\small\textbf{Definition (Complete Metric Space)}},
  fonttitle=\small\bfseries]
\begin{definition}[Complete Metric Space]
\label{def:complete-metric-space}
A metric space $(X,d)$ is \emph{complete} (or \emph{Cauchy-complete})
if every Cauchy sequence in $X$ converges to a point of $X$.
\end{definition}
\end{tcolorbox}

\begin{remark*}[Fully quantified form]
\begin{align*}
(X,d)\ \text{complete}
&\;\iff\;
\forall (x_n)_{n\in\mathbb{N}} \subseteq X, \\
&\qquad\quad
\Bigl[
  \forall \varepsilon > 0,\;
  \exists N \in \mathbb{N},\;
  \forall m,n \ge N:\;
  d(x_m,x_n) < \varepsilon
\Bigr] \\
&\qquad\Rightarrow
\Bigl[
  \exists x \in X,\;
  \forall \varepsilon > 0,\;
  \exists N' \in \mathbb{N},\;
  \forall n \ge N':\;
  d(x_n,x) < \varepsilon
\Bigr].
\end{align*}
\end{remark*}

\begin{remark*}[Why order cannot be used here]
In $\mathbb{R}$, completeness is equivalent to the least upper bound
property, which requires an ordering.
General metric spaces (e.g.\ $\mathbb{R}^2$, spaces of functions) need
not carry any compatible ordering, so the order-based formulation
cannot be transported.
The Cauchy condition provides an entirely intrinsic notion of
completeness depending only on the metric.
\end{remark*}

\begin{remark*}[Completeness is metric-dependent, not topological]
Completeness is a property of the \emph{metric}, not of the underlying
topology.
The open interval $(-1,1)$ is homeomorphic to $\mathbb{R}$ via the map
$x \mapsto \tan(\pi x/2)$, yet $\mathbb{R}$ is complete while
$(-1,1)$ with the inherited absolute-value metric is not: the sequence
$x_n = 1 - 1/n$ is Cauchy in $(-1,1)$ but converges to $1 \notin (-1,1)$.

Changing the metric on the same set can change completeness.
\end{remark*}

% ---------------------------------------------------------
\subsection{Examples}
\label{sec:completeness-examples}
% ---------------------------------------------------------

\begin{example}[Complete: $\mathbb{R}$ and $\mathbb{R}^n$]
$\mathbb{R}$ is complete: this is equivalent to the least upper bound
property and is proved in Volume II.
$\mathbb{R}^n$ with the Euclidean metric is complete by reduction to
$\mathbb{R}$: if $(x^{(m)})$ is Cauchy in $\mathbb{R}^n$, each
component sequence $(x^{(m)}_k)$ is Cauchy in $\mathbb{R}$ (since
$|x^{(m)}_k - x^{(\ell)}_k| \le d(x^{(m)}, x^{(\ell)})$), hence
converges, and the component-wise limits assemble to the limit in
$\mathbb{R}^n$.
\end{example}

\begin{example}[Complete: $C([a,b],\mathbb{R})$ with the uniform metric]
The space of continuous functions on a closed bounded interval, equipped
with
\[
  d(f,g) := \sup_{x \in [a,b]} |f(x) - g(x)|,
\]
is complete.
A Cauchy sequence $(f_n)$ in this metric converges uniformly, and the
uniform limit of continuous functions is continuous.
\end{example}

\begin{example}[Complete: any discrete metric space]
In any metric space with the discrete metric, a Cauchy sequence is
eventually constant
(see \texttt{notes-cauchy-sequences.tex}, Example),
and an eventually constant sequence trivially converges to its eventual
value.
\end{example}

\begin{example}[Incomplete: $\mathbb{Q}$]
$\mathbb{Q}$ with the standard metric is not complete.
The sequence of rational truncations of $\sqrt{2}$ is Cauchy in
$\mathbb{Q}$, but its limit $\sqrt{2}$ is irrational, hence not in
$\mathbb{Q}$.
\end{example}

\begin{example}[Incomplete: open interval $(0,1]$]
The sequence $x_n = 1/n$ is Cauchy in $(0,1]$ (it is convergent in
$\mathbb{R}$), but its limit $0$ does not belong to $(0,1]$.
\end{example}

\begin{example}[Incomplete: $C^1([0,1],\mathbb{R})$ under the uniform metric]
Smooth functions can converge uniformly to a continuous but
non-differentiable function.
For instance, $f_n(x) = \sqrt{x^2 + 1/n}$ converges uniformly to
$|x|$ on $[-1,1]$, and $|x|$ is not differentiable at $0$.
So $C^1([0,1],\mathbb{R})$ is not closed in $C([0,1],\mathbb{R})$ under the
uniform metric, and is therefore incomplete.
\end{example}

% ---------------------------------------------------------
\subsection{Completeness and the Closed-Subspace Criterion}
\label{sec:completeness-closed}
% ---------------------------------------------------------

\begin{proposition}[\texorpdfstring{\hyperref[prf:closed-subspace-complete]{Closed Subspace of a Complete Space is Complete}}{Closed Subspace of a Complete Space is Complete}]
\label{prop:closed-subspace-complete}
Let $(X,d)$ be a complete metric space and let $F \subseteq X$ be closed.
Then $(F, d|_F)$ is a complete metric space.
\end{proposition}

\begin{proof}
Let $(x_n)$ be a Cauchy sequence in $F$.
Since $F \subseteq X$ and $d|_F$ agrees with $d$ on $F$, the sequence is
also Cauchy in $X$.
By completeness of $X$, there exists $x \in X$ with $x_n \to x$.
Since $F$ is closed, the sequential characterization
(\hyperref[sec:sequential-characterizations]{\S\ref*{sec:sequential-characterizations}})
gives $x \in F$.
Hence $(x_n)$ converges to a point of $F$, so $(F, d|_F)$ is complete.
\end{proof}

\begin{proposition}[\texorpdfstring{\hyperref[prf:complete-subspace-closed]{Complete Subspace is Closed}}{Complete Subspace is Closed}]
\label{prop:complete-subspace-closed}
Let $(X,d)$ be any metric space and let $F \subseteq X$.
If $(F, d|_F)$ is complete, then $F$ is closed in $X$.
\end{proposition}

\begin{proof}
Let $(x_n)$ be a sequence in $F$ with $x_n \to x$ in $X$.
Every convergent sequence is Cauchy
(\hyperref[prop:convergent-implies-cauchy]{Proposition~\ref*{prop:convergent-implies-cauchy}}),
so $(x_n)$ is Cauchy in $F$.
By completeness of $F$, there exists $y \in F$ with $x_n \to y$ in $F$,
hence in $X$.
Uniqueness of limits in $X$ gives $x = y \in F$.
Since every sequence in $F$ converging in $X$ has its limit in $F$,
$F$ is closed.
\end{proof}

\begin{remark*}[The closed--complete equivalence]
Combining the two propositions in the setting of a complete ambient space:
\[
  \boxed{
    F\ \text{closed in complete } (X,d)
    \;\iff\;
    (F, d|_F)\ \text{is complete.}
  }
\]
This collapses two conditions into one.
In practice it means: to certify that a subspace is complete, exhibit
it as a closed subset of a known complete space (e.g.\ a closed subset
of $\mathbb{R}^n$).
To certify incompleteness, find a Cauchy sequence in the subspace with
a limit outside it --- which witnesses that the subspace is not closed.
\end{remark*}

% ---------------------------------------------------------
\subsection{Completion of a Metric Space}
\label{sec:completion}
% ---------------------------------------------------------

An incomplete metric space has ``missing'' limit points: Cauchy sequences
that converge to points not present in the space.
The completion fills in exactly those gaps and nothing more.

\begin{tcolorbox}[colback=propbox, colframe=propborder, arc=2pt,
  left=6pt, right=6pt, top=4pt, bottom=4pt,
  title={\small\textbf{Definition (Completion)}},
  fonttitle=\small\bfseries]
\begin{definition}[Completion]
\label{def:completion}
Let $(X,d)$ be a metric space.
A \emph{completion} of $(X,d)$ is a complete metric space
$(\widetilde{X}, \widetilde{d})$ together with an isometric embedding
$\iota : X \hookrightarrow \widetilde{X}$ such that $\iota(X)$ is dense
in $\widetilde{X}$.
\end{definition}
\end{tcolorbox}

\begin{remark*}[The three requirements unpacked]
\begin{enumerate}[label=(\roman*), leftmargin=2em]
  \item \textbf{Isometric:}
        $\widetilde{d}(\iota(x), \iota(y)) = d(x,y)$ for all $x,y \in X$.
        The embedding preserves all distances; no geometry is distorted.
  \item \textbf{Dense:}
        $\overline{\iota(X)} = \widetilde{X}$.
        Every point of $\widetilde{X}$ is a limit of a sequence from
        $\iota(X)$.
        The completion adds only what is forced by limits already present
        in $X$.
  \item \textbf{Complete:}
        $(\widetilde{X}, \widetilde{d})$ itself has no gaps.
\end{enumerate}
\end{remark*}

\begin{remark*}[The co-Cauchy construction]
\label{rem:co-cauchy}
Two Cauchy sequences $(p_n)$ and $(q_n)$ in $X$ are \emph{co-Cauchy} if
$d(p_n, q_n) \to 0$.
Co-Cauchyness is an equivalence relation on the set of all Cauchy sequences
in $X$.

The completion $\widetilde{X}$ is defined as the collection of
equivalence classes $[(p_n)]$, with metric
\[
  \widetilde{d}\bigl([(p_n)],\, [(q_n)]\bigr)
  \;:=\;
  \lim_{n\to\infty} d(p_n, q_n).
\]
That this limit exists follows from the four-point lemma:
$|d(p_m,q_m) - d(p_n,q_n)| \le d(p_m,p_n) + d(q_m,q_n)$,
which shows $(d(p_n,q_n))$ is Cauchy in $\mathbb{R}$, hence convergent.

Each $x \in X$ embeds as the constant sequence $\iota(x) = [(x,x,x,\ldots)]$,
and this embedding is isometric.
\end{remark*}

\begin{tcolorbox}[colback=thmbox, colframe=thmborder, arc=2pt,
  left=6pt, right=6pt, top=4pt, bottom=4pt,
  title={\small\textbf{Theorem (Completion Theorem)}},
  fonttitle=\small\bfseries]
\begin{theorem}[\texorpdfstring{\hyperref[prf:completion]{Completion Theorem}}{Completion Theorem}]
\label{thm:completion}
Every metric space $(X,d)$ has a completion.
The completion is unique up to isometric isomorphism fixing $X$.
\end{theorem}
\end{tcolorbox}

\begin{remark*}[Proof outline]
The co-Cauchy construction gives existence in three steps.
\begin{enumerate}[label=(\roman*), leftmargin=2em]
  \item \textbf{$\widetilde{d}$ is a well-defined metric.}
        The four-point lemma guarantees the limit defining $\widetilde{d}$
        exists and is independent of the choice of representatives.
  \item \textbf{$\iota(X)$ is dense in $\widetilde{X}$.}
        The equivalence class $[(p_n)]$ is the $\widetilde{d}$-limit of
        the sequence of constant classes $[(\iota(p_n))]$.
  \item \textbf{$\widetilde{X}$ is complete.}
        Given a Cauchy sequence $(P_k)$ of equivalence classes, one
        selects representatives with terms mutually close, extracts a
        diagonal sequence $(q_n)$ that is Cauchy in $X$, and shows
        $[(q_n)]$ is the $\widetilde{d}$-limit of $(P_k)$.
\end{enumerate}
Uniqueness is proved by showing any two completions are connected by a
unique isometric isomorphism fixing the image of $X$.
\end{remark*}

\begin{remark*}[Canonical example: $\mathbb{Q} \leadsto \mathbb{R}$]
The completion of $\mathbb{Q}$ under the standard metric is isometrically
isomorphic to $\mathbb{R}$.
This gives an alternative construction of the real numbers from rational
Cauchy sequences, without Dedekind cuts.
Arithmetic on $\mathbb{R}$ is defined class-wise:
$[(p_n)] + [(q_n)] = [(p_n + q_n)]$, and so on.
\end{remark*}

\begin{remark*}[Source note]
The Completion Theorem is Theorem 2.80 in Pugh (\S10).
The co-Cauchy construction follows Pugh's proof closely.
An alternative embedding argument (embedding $X$ isometrically into the
complete space $C^0_b(X,\mathbb{R})$ of bounded continuous functions)
is sketched in Pugh Exercise 4.39.
\end{remark*}

% ---------------------------------------------------------
\subsection{Conceptual Summary}
\label{sec:completeness-summary}
% ---------------------------------------------------------

\begin{center}
\begin{tikzpicture}[
  >=Latex,
  node distance=2.6cm,
  box/.style={draw, rounded corners=3pt, align=center, inner sep=8pt},
  lab/.style={font=\small},
  impl/.style={->, thick},
  noimpl/.style={->, thick, dashed}
]

\node[box] (conv)     {Convergent\\$(x_n)\to x\in X$};
\node[box, right=3.8cm of conv] (cauchy)   {Cauchy\\$d(x_m,x_n)\to 0$};
\node[box, below=2.2cm of cauchy]           (bounded)  {Bounded};
\node[box, below=2.2cm of conv]             (complete) {$X$ complete};

\draw[impl]   (conv)    -- node[lab, above]{always} (cauchy);
\draw[noimpl] (cauchy)  -- node[lab, above]{only if $X$ complete} (conv);
\draw[impl]   (cauchy)  -- node[lab, right]{always} (bounded);
\draw[impl]   (cauchy)  edge[bend right=20]
              node[lab, left]{in complete $X$} (conv);
\draw[impl]   (complete.north) -- node[lab, right, xshift=3pt]{enables} (conv.south);

\end{tikzpicture}
\end{center}

\begin{center}
\[
  \boxed{
    \text{Complete} = \text{``every Cauchy sequence converges''}
  }
\]
\[
  \boxed{
    \text{Every metric space embeds isometrically into a complete one.}
  }
\]
\[
  \boxed{
    F\ \text{closed in complete } X
    \;\iff\;
    F\ \text{itself complete.}
  }
\]
\end{center}

\begin{remark*}[Forward reference]
Completeness interacts with compactness: every compact metric space is
complete, but the converse fails (e.g.\ $\mathbb{R}$ is complete but not
compact).
In $\mathbb{R}^n$, Heine--Borel identifies compact sets as precisely the
closed and bounded ones, linking completeness, closure, and boundedness.
These connections are developed in \texttt{notes-compactness.tex}.
\end{remark*}